\section{Open challenges and a research roadmap}
\label{sec:challenges}

\paragraph{R1: The ontology construction bottleneck.} No public domain
ontology covers a bottling line, a paper machine or a battery cell assembly at
the granularity an agent needs, and hand-building one per plant does not scale.
The promising direction is bootstrapping from artefacts that already encode the
structure, \opcua{} address spaces, AutomationML instance hierarchies, P\&IDs
via DEXPI, PLC projects via PLCopen~XML, with an LLM proposing the alignment
and a formal checker accepting or rejecting it. The research question is not
whether models can draft ontologies (they can) but how to obtain a
\emph{soundness guarantee} on the result; our lift-validate-repair pattern is
one answer, and its limits are unexplored.

\paragraph{R2: Publishing constraints, not just data.} Interlocks, operating
envelopes and safety conditions are the highest-value semantic content in a plant
and the least likely to be machine-readable. We see no technical obstacle to a
standardised \emph{constraint submodel}, SHACL shapes or an equivalent, carried
alongside an AAS or an \opcua{} companion model, versioned and signed. This is,
in our view, the single most useful standardisation action the community could
take for agent safety.

\paragraph{R3: Temporal and streaming semantics.} We modelled observations as
points. Real industrial reasoning is about trends, windows, rates of change and
episodes. RDF-star, named graphs and stream-reasoning proposals exist but no
convention has taken hold, and the token economics of putting time series in a
context window are hopeless. The likely answer is the E7 pattern generalised:
temporal aggregation as a tool, with the semantic layer describing what the
series \emph{is} and the engine computing over it.

\paragraph{R4: Text-to-query accuracy over industrial schemas.} The query
architecture we recommend transfers the accuracy problem to query generation.
Industrial graphs are unlike the encyclopaedic graphs on which text-to-SPARQL is
usually evaluated: schemas are deeper, identifiers are opaque, and errors are
consequential. A benchmark for query generation over standardised industrial
models, with execution-based scoring and a distinction between wrong and
unsafe, does not exist and should.

\paragraph{R5: Authorisation semantics for autonomous agents.} An agent that
writes needs the OT analogue of an operator's competency and shift authorisation:
what may this agent do, to which assets, under which plant states, with what
human confirmation. ODRL supplies the expression language and IEC~62443
\cite{iec62443} the security framing, but nothing joins them to the equipment
model. The
lift-and-validate pattern extends naturally, validate the proposed action
against the policy graph as well as the plant graph, and the composition is
unstudied.

\paragraph{R6: Benchmarks with actions and consequences.} Existing grounding
and hallucination benchmarks evaluate answers. Industrial agents emit actions,
and the interesting metrics are asymmetric: a missed detection and a false alarm
have very different costs, and both differ by fault family. The corpus released
with this paper is a starting point, twelve labelled fault families with a
valid control set, but it is one plant and synthetic. A community benchmark
built from anonymised real installations, with severity-weighted scoring, is
needed.

\paragraph{R7: Semantic model verification as a first-class problem.} If
agents are to depend on the semantic layer, the semantic layer must itself be
verified: are the units on these two tags consistent with the equation that
relates them, does the declared topology match the P\&ID, is the interlock model
equivalent to the ladder logic it claims to represent? Each of these is a
decidable check, none is routine practice, and an incorrect semantic layer is
more dangerous than none because it is trusted.

\paragraph{R8: Learned projection policies.} Finding~F10 shows that deciding
\emph{which semantic layer a question needs} is worth more than enlarging the
context window, and that a lexicon of surface cues already recovers most of the
gap. A learned router, trained on the fact kinds a question turns out to
require, which the benchmark supplies as labels, should do better, and would
turn context assembly into an optimisation problem with a measurable objective
(coverage per token) rather than a heuristic. The same machinery would decide
when to stop projecting and issue a query instead, which is currently a manual
design decision.

\paragraph{R9: Retiring convention-carried semantics.} Finding~F9 shows that
a substantial share of an un-grounded system's competence rests on tag
mnemonics, and that the same inference becomes an active liability in plants
where the convention has decayed. Detecting that decay is tractable and, as far
as we know, unattempted: an agreement check between what a tag name asserts and
what the information model states would flag exactly the tags on which a
language model will be confidently wrong. Every plant with a hand-maintained
cross-reference spreadsheet already knows it needs this.

\paragraph{R10: Measuring the operator.} Finding~F4, that weak validation
over-blocks, has a human consequence we did not measure. The rate at which an
assistant refuses legal instructions is plausibly a stronger determinant of
adoption than its error rate, and it is directly controlled by the quality of the
semantic layer. This connects semantic engineering to trust calibration, and is
straightforwardly testable.
